\chapter*{Abstract}
\addcontentsline{toc}{chapter}{Abstract}

The rapid advancement of Artificial Intelligence (AI) and Large Language Models (LLM) has created new opportunities for the enhancement of educational support systems. However, many existing AI tools generate generic explanations without considering the structured nature of university syllabi, often resulting in incomplete or unbalanced study materials. Students face challenges in converting syllabus topics and reference materials into organized, exam-oriented notes aligned with academic requirements frequently. This project proposes the development of an AI-based academic support system capable of generating syllabus-aligned study notes using a retrieval-augmented large language model architecture. The system processes academic documents such as syllabus outlines and reference materials, converts them into vector embeddings, and stores them in a vector database to enable efficient retrieval of relevant content. When a user requests notes for a particular chapter or topic, the system retrieves the most relevant information and generates structured academic notes. To ensure balanced coverage and alignment with syllabus requirements, the proposed system generates detailed notes for each topic based on the retrieved reference materials and the syllabus context. The output will be presented in a structured Markdown format optimized for exam preparation. The proposed system aims to assist students in efficiently transforming syllabus topics into organized learning materials while ensuring comprehensive topic coverage and clarity. By integrating retrieval-based knowledge systems with large language models, the project demonstrates how AI can be applied to improve academic learning tools and support effective examination preparation. \par

\vspace{1cm}

\textit{\textbf{Keywords:} Artificial Intelligence, Educational Support System, Large Language Model, Retrieval-augmented }